\documentclass[10pt,a4paper]{amsart}
\usepackage[margin=19mm]{geometry}
\usepackage{amsmath,amssymb,mathtools}
\usepackage[scr=rsfs,cal=euler]{mathalfa}
\usepackage{xfrac}
\usepackage[unicode,hidelinks,bookmarks=false]{hyperref}
\newcommand{\ie}{i.e.\ }
\newcommand{\cf}{cf.\ }
\newcommand{\resp}{respectively\ }
\newcommand{\Subsection}{\subsection}
\providecommand{\nobreakdash}{\nobreak\hbox{-}}
\setlength{\emergencystretch}{2em}
\multlinegap0pt
\usepackage{amssymb}
\usepackage{mathtools}
\usepackage[shortlabels]{enumitem}
\usepackage{xcolor}
\newtheorem{proposition}{Proposition}
\title{Equivariant irrationality of very general symmetric Verra fourfolds}
\author{Aideen Fay}
\date{Source: arXiv:2605.30439v1 (CC BY 4.0)}
\begin{document}
\maketitle

\noindent\emph{Source and licence.} Equivariant irrationality of very general symmetric Verra fourfolds. Version arXiv:2605.30439v1 (CC BY 4.0), distributed by arXiv under the Creative Commons Attribution 4.0 International licence, http://creativecommons.org/licenses/by/4.0/, as recorded in the arXiv licence metadata for this exact version and shown on the abstract page https://arxiv.org/abs/2605.30439v1. Full licence notice: \texttt{licenses/arxiv-2605.30439-CC-BY-4.0.txt}.\par\smallskip

\noindent\textit{Editorial scope.} Counted unit: the article's proposition proposition (source0.tex lines 1050--1059). The article's complete original proof of it is delivered with it (source0.tex lines 1061--1081, one \texttt{proof} environment of 21 lines). No mathematics is written, repaired or re-derived: the statement and the proof are the article's own lines, in the article's own notation, and no retained line is substituted or re-flowed.  No further source-local context is needed: every cross-reference of every retained line resolves inside the two delivered spans.  Nothing was omitted from any retained span, so the package carries no omission record.  No retained line contains a citation, so the wrapper carries no bibliography and no bibliography entry.  No retained line contains a figure environment, an image-inclusion command or drawing code, so no image is delivered and the package declares no asset dependency.  Editorial formatting only, in addition: an AMS \texttt{amsart} preamble that restates the article's own declarations and macros used by the retained lines, namely the environments \texttt{proposition} on separate counters, together with the standard packages those lines need.  The retained environments print in this document as Proposition 1 in order of appearance, whereas the article prints the same items with the numbering of its own full document.  The complete native source of the article as distributed in the arXiv e-print for this exact version is bundled unchanged as \texttt{sources/201674/source0.tex}.
\begin{proposition}
\label[proposition]{lemma:quartic}
Let $B=\{f_{2,2}=0\}\subset \mathbb P^2\times\mathbb P^2$
be the smooth $(2,2)$-divisor and let $\Delta_{\mathrm{diag}}\subset \mathbb P^2\times\mathbb P^2$
be the diagonal. Then
\[
  C_4:=B\cap \Delta_{\mathrm{diag}}
\]
is a smooth plane quartic.
\end{proposition}

\begin{proof}
Suppose \(p\in C_4\) were singular. Then
\[
  d(f_{2,2}|_{\Delta_{\mathrm{diag}}})_p=0.
\]
By symmetry of \(f_{2,2}\), at the point \((p,p)\in\Delta_{\mathrm{diag}}\) we have
\[
  df_{2,2}|_{(p,p)}=(\alpha,\alpha)
  \in T_p^\vee\mathbb P^2\oplus T_p^\vee\mathbb P^2.
\]
Restricting to the diagonal gives
\[
  d(f_{2,2}|_{\Delta_{\mathrm{diag}}})_p=2\alpha.
\]
Since we work over \(\mathbb C\), this implies \(\alpha=0\). Hence
\[
  df_{2,2}|_{(p,p)}=0.
\]
As \(p\in C_4\), we also have \(f_{2,2}(p,p)=0\). Therefore \(B\) would be singular
at \((p,p)\), which contradicts the smoothness of \(B\). Thus \(C_4\) is smooth.
\end{proof}

\end{document}
